\section{Introduction}
\label{sec:intro}

A recurrent or looped model applies the same block again and again, and every extra iteration buys more
computation on the same input. This lets small networks extrapolate to instances harder than any seen
in training, provided they are allowed to think for longer
\citep{schwarzschild2021thinking,bansal2022endtoend,geiping2025scaling,jolicoeur2025trm}. The promise is
adaptive compute, with few steps on easy inputs and many on hard ones, and keeping it requires a
stopping rule. The rules in common use, learned halting gates \citep{graves2016act,banino2021pondernet},
confidence thresholds and convergence tests \citep{zhou2020pabee,geiping2025scaling}, are all fitted on
the difficulty seen in training. Yet the inputs that matter most are the ones harder than anything in
the training set, and on them the rule must decide that the computation is finished from information it
was never rewarded for reading.

We study this decision in a setting where the right answer is known. In a \emph{local} recurrent
reasoner, one step moves information by at most $r$ positions, so after $t$ steps position $i$ depends
only on the input within radius $rt$, its light cone. If that window does not yet determine the target,
some other input agrees on the window but has a different answer, and the model must be wrong on one of
the two. Every input therefore has a lower bound $\LB(x)=\ceil{\dstar(x)/r}$ on the first step at which
its answer can be certified (\cref{thm:lb}), where the \emph{depth} $\dstar(x)$ is the largest radius any
position needs. Halting earlier returns an answer that the model's own state does not determine, correct
only by luck of the data. Because the bound depends on depth and not on length, it separates two shifts
that are usually confounded, longer inputs and deeper ones. For the tasks we study, a constant-size local
automaton halts at exactly $\LB(x)$ at every length (\cref{prop:automaton}), so the bound is attainable
and we can ask which stopping rules attain it.

Our contributions are threefold. We prove the light-cone bound, which turns ``halting too early'' into a
checkable per-instance quantity, derive its consequences for local commit-on-done rules and for rules
that see only the length, and give closed-form depth oracles and a matching automaton for addition and
segmented parity, each with an executable test (\cref{sec:theory}). We introduce DepthShift, whose
frozen splits vary depth and length independently up to $12.8\times$ the training depth
(\cref{sec:benchmark}). And we evaluate fifteen stopping rules and four references on one backbone family,
from post-hoc confidence and convergence tests through ACT, PonderNet and pooled TRM-style and v1 heads to
a family of local \emph{completion heads} that we introduce, scoring each by accuracy, steps, and how often
it stops or errs before $\LB(x)$ (\cref{sec:method,sec:experiments}).

The results tell a consistent story. Jointly trained pooled gates leak length: on addition, ACT is
99.8\% accurate on deeper inputs with fewer steps than the bound allows, because it reads a global
shortcut, and falls to 35.9\% once inputs are both deeper and longer. Premature stopping and premature
error are different things: convergence rules halt before the bound on most deep addition inputs and
are still right, because a travelling carry keeps the hidden state moving until it lands. When pending
positions are genuinely ambiguous, as in addition, per-position confidence is already a certificate;
max-softmax halts at exactly $\LB(x)$ on every split, and our completion heads match it without improving
on it. When they are not, as in parity, the completion heads learn the training clock, while pooled gates
extrapolate in depth but fall below 1\% once inputs are also longer. Finally, jointly trained completion
heads leave the backbone's solution intact, whereas gates that reweight the task loss degrade it. The
difficulty of termination detection thus lies less in the readout than in training data that never show
a computation still unfinished at a late step.
